%% main.tex — E-Health for All Manuscript
%% Converted from Group10_Manuscript_FINAL_Pre_Oral.docx
%% Compile: pdflatex main.tex → bibtex main → pdflatex main.tex × 2

\documentclass{research}

%% PDF metadata
\hypersetup{
  pdftitle  = {E-Health for All: A User-Centered Telemedicine System},
  pdfauthor = {Bruzo, Ian; Narvato, Dominnica SL.; Paraiso, Kristine Carla I.; Villamer, EJ Andrei J.}
}

%% ── Citation key map (document order → bib key) ────────────────────────────
%% The manuscript uses numeric citations [1]–[72].
%% We use the cite package for IEEE-style compressed brackets.

\begin{document}

%% ═══════════════════════════════════════════════════════════════════════════
%%  TITLE PAGE
%% ═══════════════════════════════════════════════════════════════════════════
\maketitlepage
  {\textit{E-Health for All}: A User-Centered Telemedicine System}
  {An Undergraduate Research Paper}
  {College of Computer Studies \\ Naga College Foundation, Inc. \\ Naga City}
  {In Partial Fulfillment of the Requirements for the Degree}
  {Bruzo, Ian \\[4pt]
   Narvato, Dominnica SL. \\[4pt]
   Paraiso, Kristine Carla I. \\[4pt]
   Villamer, EJ Andrei J.}
  {2027}
  {Bachelor of Science in Information Systems}

%% ═══════════════════════════════════════════════════════════════════════════
%%  SECTION 1 — INTRODUCTION
%% ═══════════════════════════════════════════════════════════════════════════
\setcounter{page}{1}
\section{Introduction}

Healthcare access remains one of the most persistent global inequalities, with limited specialist
availability and inadequate infrastructure keeping millions from critical care. The World Health
Organization estimates that half of the global population lacks access to essential health services,
with rural communities disproportionately affected \cite{WHO2023}, leading to delayed diagnoses and
preventable deaths. Telemedicine presents a promising solution by delivering healthcare across
distances, connecting remote patients with specialists without requiring travel and making scarce
resources available to underserved populations.

In the Philippines, these challenges are especially severe. The country's geography and urban
concentration of medical specialists leave rural populations with limited healthcare options. In
response, the Philippine government and health institutions have begun embracing telemedicine to
broaden coverage. The National Telehealth Center has launched SMS and web-based platforms
connecting rural patients with specialist consultations \cite{Macariola2021,Pickering2024}, and
various health institutions have adopted electronic medical records to improve continuity of care
across facilities \cite{Santos2021,PhilPrimaryStudies2024}.

Despite these efforts, significant barriers persist. Poor transportation infrastructure hinders
facility access, particularly during emergencies and follow-up care. Unreliable internet
connectivity, limited computing resources, and low digital literacy among both patients and
healthcare workers further compound these challenges \cite{Abdul2024,Castillo2023}. Existing
telemedicine initiatives, while promising, remain fragmented and rarely address the unique
technological constraints and cultural context of rural Filipino communities
\cite{Domingo2023,Cruz2022}.

This research seeks to bridge these gaps by developing an integrated telemedicine platform
specifically designed for the selected rural barangays. The platform combines three components—a
comprehensive doctors' database, a patient records management system, and an appointment
scheduling system—into one user-friendly interface accessible to users with different levels of
digital literacy.

This initiative aligns with multiple UN Sustainable Development Goals: good health and well-being
(SDG 3), by improving healthcare access and quality. It supports industry, innovation, and
infrastructure (SDG 9), by demonstrating how technology can address healthcare infrastructure gaps
in underserved areas. It advances reduced inequalities (SDG 10), by examining how digital health
solutions reduce disparities in healthcare access across different populations and socioeconomic
groups.

\subsection*{Telemedicine Adoption and Implementation}

Telemedicine, the delivery of healthcare services using information and communication
technologies, has emerged as a vital solution for healthcare access issues in isolated and
underserved communities. Telemedicine allows healthcare providers to assess, diagnose, and treat
patients in remote areas by utilizing telecommunications technology to effectively connect medical
experts with populations that have limited access to healthcare facilities. It has shown great
promise in overcoming healthcare access issues in rural communities around the world.

Research shows that telemedicine consultations promote better clinical decision-making, reduce
unnecessary patient transfers, and increase the chances of local admissions instead of costly
transfers to far-off hospitals \cite{Natafgi2021,McCormick2021}. These advantages apply to many
medical fields, from primary care to specialized areas like neurology and emergency medicine
\cite{Ezeamii2024}.

However, the pandemic highlighted the shortcomings of telemedicine: access services at lower
rates \cite{Chu2021,Kichloo2020}. Digital literacy among older adults and low-income individuals
\cite{Abdul2024,Castillo2023}, equal access concerns \cite{Fairchild2019,DAmore2021,Anzalone2025},
implementation costs and maintenance \cite{Zachrison2021,Chen2020}, infrastructure issues
\cite{Chandrasekaran2021,Ngugi2021,Rhoades2025}, and data security and patient privacy concerns
\cite{Smith2021}.

On the one hand, communities with integrated systems achieve better health outcomes for managing
chronic diseases compared to those that use fragmented record-keeping \cite{Mendoza2023}. Patient
satisfaction with telemedicine in the Philippines shows encouraging trends \cite{Mendoza2023,Salazar2021}.
Even senior citizens increasingly accepted remote consultations \cite{Aquino2022}. The University
of the Philippines Health Service successfully used tailored electronic medical records systems in
various settings during COVID-19, ensuring continuity of patient care despite mobility restrictions
\cite{PhilPrimaryStudies2024,Pickering2024}. Thorough assessments show considerable progress in
adoption over time, although significant gaps between different types of facilities and locations
persist \cite{AdlerMilstein2023,Rhoades2025}.

Successfully implementing telemedicine requires careful adjustments to local circumstances.
Detailed analyses reveal Philippines-specific barriers like governance issues, unclear licensure
requirements, unresolved liability problems, and vague reimbursement mechanisms that hinder
healthcare provider involvement \cite{Macariola2021,Domingo2023}. Healthcare technology must
consider local cultural practices, language preferences, and health-seeking behaviors to ensure
meaningful user acceptance in diverse Filipino communities \cite{Cruz2022}. Creating user-friendly
systems is crucial for sustained adoption rather than leaving behind abandoned pilot projects
\cite{Tiu2022}.

\subsection*{Patient Information Management System}

Patient information management involves the organized collection, storage, and use of health data
to support clinical decision-making and ensure continuity of care. Electronic Medical Records (EMR)
systems digitize patient health information, enabling providers to access complete histories, track
treatments, identify medication interactions, and coordinate care across specialties. This is
especially critical in rural and underserved areas where access to prior care records is limited.

EMR systems significantly improve care coordination, with integrated communities achieving better
chronic disease outcomes compared to those with fragmented records \cite{Mendoza2023,Gill2020}.
However, while basic adoption rates are high, comprehensive implementation remains lower in rural
facilities despite recognized benefits \cite{DAmore2021,Anzalone2025}. Technical, organizational,
and workflow barriers continue to hinder smooth information sharing \cite{Ngugi2021,AdlerMilstein2023,Rhoades2025},
and digital literacy gaps among older adults and low-income individuals require targeted design and
training interventions \cite{Abdul2024}.

Data security and patient privacy are foundational to EMR trust, requiring encryption, access
controls, and audit trails \cite{Smith2021}. Integrated systems reduce administrative burdens and
improve care delivery \cite{Betancor2024,Kucuk2024,Mazaheri2024}, while predictive models using
patient data enable targeted interventions and better resource allocation \cite{Ferro2020,Dantas2018},
supporting a shift toward proactive, preventive primary care \cite{Cusack2020}.

In the Philippines, EMR adoption shows both progress and challenges. COVID-19 visibly impacted
system usage in rural health facilities \cite{Acacio2024}, yet implementation improved clinical
decision-making and reduced medication errors compared to paper systems \cite{Santos2021}.
Sustained adoption is linked to training, technical support, and gradual workflow integration
rather than rapid rollouts \cite{Mendoza2023}, though research representation across specialties,
settings, and facility types remains uneven \cite{Derecho2024}. The absence of centralized
provider databases causes referral delays and disjointed care \cite{Reyes2023}, while comprehensive
physician databases with specialization and scheduling details enhance diagnostic workflows
\cite{Hernandez2022}. Keeping such databases accurate remains challenging, necessitating automated
verification and validation processes \cite{Patel2023}.

\subsection*{Appointment Scheduling in Healthcare}

Appointment scheduling systems are essential to healthcare delivery, coordinating patient needs,
provider availability, and facility resources. Well-designed systems reduce waiting times, lower
missed appointments through automated reminders, improve provider productivity, and ensure urgent
cases receive prompt attention. In rural settings, where patients face long travel distances and
transportation barriers, efficient scheduling is especially critical for maintaining access and
treatment adherence.

Digital scheduling significantly reduces no-show rates compared to offline booking, with automated
reminders proving particularly effective \cite{Betancor2024}. AI-driven tools further improve
efficiency by predicting non-attendance and optimizing resource allocation \cite{Kucuk2024}, while
open-access scheduling reduces no-shows by offering same-day or next-day appointments
\cite{Mazaheri2024}. Predictive models incorporating patient demographics and appointment history
identify high lead time and prior missed appointments as the strongest predictors of non-attendance
\cite{Ferro2020,Dantas2018}. Integrated systems also improve patient flow and reduce waiting area
congestion in underserved facilities \cite{Wang2023}, supporting broader proactive primary care
initiatives through timely referrals and fewer missed appointments \cite{Cusack2020,Patel2023}.

In the Philippines, digital appointment management in rural health centers has shortened waiting
times and reduced missed appointments through better time slot management and automated reminders
\cite{DelRosario2024}. However, rural contexts require flexible approaches, as rigid scheduling is
ineffective where patients face unreliable transportation or inflexible work schedules. Systems
must therefore incorporate easy rescheduling options and reminder methods that function without
internet access \cite{Ramirez2021}.

\subsection*{Synthesis of the State-of-the-Art}

The detailed collection of works and studies offered valuable insights into the topic. These studies
backed their claims with relevant data and detailed information. This section provides a brief
overview of the reviews of relevant literature.

Telemedicine has emerged as a viable solution to geographic barriers, improving clinical
decision-making, reducing unnecessary transfers, and expanding specialist access
\cite{Ezeamii2024,Natafgi2021,McCormick2021}. However, persistent challenges remain, including
infrastructure limitations, digital literacy gaps, data security concerns, and Philippines-specific
barriers such as unclear licensure, liability, and reimbursement mechanisms
\cite{Macariola2021,Smith2021,Castillo2023,Domingo2023}. Despite these, large-scale pilots in the
Philippines demonstrate feasibility across diverse communities, with patient satisfaction showing
encouraging trends even among senior citizens \cite{Aquino2022,Noceda2023,Salazar2021,Pickering2024}.
Successful adoption requires culturally sensitive, user-centered design tailored to local contexts
\cite{Cruz2022,Tiu2022}.

Patient information management systems significantly improve care coordination and clinical
outcomes, yet comprehensive EMR implementation remains lower in rural facilities
\cite{Anzalone2025,DAmore2021}. Workflow barriers, digital literacy gaps, and data security
concerns continue to hinder full adoption \cite{Abdul2024,AdlerMilstein2023,Rhoades2025,Smith2021}.
In the Philippines, gradual integration supported by training and technical assistance yields
better sustained adoption \cite{Mendoza2023}, while the absence of centralized provider databases
causes referral delays and fragmented care \cite{Hernandez2022,Reyes2023}.

Appointment scheduling systems directly reduce no-show rates, improve patient flow, and support
proactive primary care \cite{Cusack2020,Betancor2024,Kucuk2024}. Predictive modeling and
AI-driven tools further optimize resource allocation \cite{Ferro2020,Dantas2018,Kucuk2024}. In
Philippine rural health centers, digital scheduling has measurably shortened waiting times
\cite{DelRosario2024}, though flexible, offline-capable systems are essential given transportation
and connectivity constraints \cite{Ramirez2021}.

Across all three areas, the literature consistently points to the need for an integrated,
context-sensitive platform that unifies patient records, provider information, and appointment
scheduling—precisely the gap this research addresses.

\subsection*{Gap Bridged by the Study}

This study addresses important gaps in rural healthcare delivery by creating an integrated
telemedicine platform specifically for rural barangays in Naga City. Existing telemedicine systems
often look at patient management and appointment scheduling separately. This research offers a more
complete solution by combining a doctors' database, a patient records management system, and
appointment scheduling into one platform. It addresses the unreliable internet connectivity that is
common in rural areas. Healthcare workers can keep entering data and accessing patient information
even without internet access. Once the connection is restored, the system syncs automatically. By
bringing these features together in one practical system designed for resource-limited settings,
this research shows how telemedicine solutions can be adapted to meet the real technological,
infrastructure, and cultural challenges faced in rural healthcare.

\subsection*{Theoretical Framework}

This study is anchored on the following theories: the Task--Technology Fit (TTF) Theory by
Goodhue and Thompson, as cited by Thanthrige et al.\ \cite{Thanthrige2025}; the Diffusion of
Innovations (DOI) Theory by Rogers, as cited by Tsogbe et al.\ \cite{Tsogbe2026}; and the
Technology Acceptance Model (TAM) by Davis, as cited by Holtz et al.\ \cite{Holtz2022}.

\begin{figure}[H]
  \centering
  % Placeholder: replace with actual figure file
  % \includegraphics[width=0.85\textwidth]{figures/theoretical_paradigm.png}
  \fbox{\parbox{0.85\textwidth}{\centering\vspace{2cm}
    \textit{[Figure 1: Theoretical Paradigm --- insert image here]}
    \vspace{2cm}}}
  \caption{Theoretical Paradigm}
  \label{fig:theoretical}
\end{figure}

The diagram illustrates the theoretical foundations guiding the development of the E-Health for
All telemedicine system for improving access in rural communities.

Goodhue and Thompson (1995) established that technology is most effectively utilized when its
capabilities align with the tasks users must perform. Thanthrige et al.\ \cite{Thanthrige2025}
reinforced this principle in a 2025 systematic review, confirming that task-technology alignment
remains a critical determinant of healthcare technology adoption. This theory anchors the study
by evaluating whether the three core features of E-Health for All---digital patient records
management, automated health worker scheduling, and online appointment booking---fit the actual
daily tasks of Panicuason health facility personnel, as a technically sound system that lacks
this fit will be underused or abandoned.

Rogers (2003) argued that technology adoption is shaped by five key attributes: relative
advantage, compatibility, complexity, trialability, and observability. Tsogbe et al.\
\cite{Tsogbe2026} applied this framework in 2026 within a telehealth adoption context in a
developing-region setting, affirming its continued relevance in guiding digital health
implementation. This theory guides the study by ensuring that the E-Health for All platform
offers a clear operational advantage over the current paper-based system while remaining
compatible with the existing workflows of Barangay Panicuason.

Davis (1989) proposed that technology adoption is primarily driven by perceived usefulness---the
degree to which a system improves work performance---and perceived ease of use---the degree to
which it is free of effort. Holtz et al.\ \cite{Holtz2022} applied TAM in 2022 specifically to
telemedicine adoption among rural populations, affirming its applicability in this context. This
theory directly shapes the acceptability evaluation instrument of the study by measuring whether
health facility personnel perceive E-Health for All as useful, easy to use, and trustworthy
enough to adopt.

\subsection*{Conceptual Framework}

The Input-Process-Output (IPO) Framework, added with Feedback loop as the fourth component, is
the approach utilized in the conceptualization of this study.

\begin{figure}[H]
  \centering
  % \includegraphics[width=0.85\textwidth]{figures/conceptual_paradigm.png}
  \fbox{\parbox{0.85\textwidth}{\centering\vspace{2cm}
    \textit{[Figure 2: Conceptual Paradigm --- insert image here]}
    \vspace{2cm}}}
  \caption{Conceptual Paradigm}
  \label{fig:conceptual}
\end{figure}

The enhanced system paradigm is the approach utilized in the conceptualization of this study. The
input came from various sources to gain a complete understanding of rural healthcare needs and
readiness for technology, such as document analysis of health records, observations of
healthcare-seeking behaviors and travel times, surveys on evaluating user satisfaction and
technology skills, and interviews gathering in-depth information about access challenges and
provider readiness for telemedicine services. The processes involved include systems analysis and
specification for system design, functional testing to ensure adherence to system requirements,
usability testing to ensure the system meets user needs, and acceptability evaluation to validate
the system feasibility for deployment. The study's output is a web application with features
including computer vision, security authentication, verification, patient check-up history, QR
code security; by using the system it is expected to improve efficiency, achieve fewer
documentation burdens, shorter waiting times through better appointment management, enhance user
satisfaction, and improve data security of patient information. The continuous feedback mechanism
will ensure iterative refinement of the system to remain user-friendly, operational, and aligned
with rural healthcare.

\subsection*{Statement of the Problem}

This study aims to address the problem in healthcare accessibility of rural populations.
Specifically, this study seeks to answer the following questions:

\begin{enumerate}
  \item What are the current practices of rural healthcare providers in rural populations with
        respect to:
        \begin{itemize}
          \item Patient records management
          \item Doctor scheduling
          \item Appointment systems
        \end{itemize}
  \item What are the challenges encountered in the current system?
        \begin{itemize}
          \item Patient records management
          \item Doctor scheduling
          \item Appointment systems
        \end{itemize}
  \item What system can be designed to address the challenges encountered?
  \item What is the level of acceptability of the system in terms of:
        \begin{itemize}
          \item Usability
          \item Reliability
          \item Data security
        \end{itemize}
\end{enumerate}

\subsection*{Assumptions}

The research study is anchored on the following key assumptions:

\begin{itemize}
  \item Rural healthcare providers in Panicuason, Naga City have specific practices for managing
        patient records, scheduling doctors, and handling appointments.
  \item There are challenges with the current systems for managing patient records, scheduling
        doctors, and setting up appointments in rural healthcare facilities.
  \item A user-centered telemedicine platform can be designed and developed to tackle these
        challenges in the current healthcare delivery system.
  \item The rural healthcare providers and clinic staff are ready to accept and shift from the
        current manual practice towards the proposed user-friendly telemedicine platform and
        digital records management system.
\end{itemize}

\subsection*{Scope and Delimitations of the Study}

This study focuses on the design and development of a user-centered telemedicine platform for
Barangay Panicuason Rural Health Unit in Naga City, intended to improve patient records
management, health worker scheduling, and appointment systems for rural healthcare delivery.

The study covers the full design and development of a web-based prototype with the following
features: digital patient health profiles, clinical encounter documentation, vital signs
monitoring, medication tracking, medical history access, appointment generation, and an integrated
doctor database. Respondents are health facility personnel of Barangay Panicuason Rural Health
Unit, consisting of physicians, nurses, midwives, and administrative staff, as well as rural
community residents who are primary users of the facility. The system is built using PHP, HTML,
CSS, and JavaScript with a MySQL database on a XAMPP local server, and is compatible with Windows
10/11, macOS, Linux, and Android devices (version 8.0 or higher, minimum resolution of
$720\times1280$ pixels).

This study is delimited to Barangay Panicuason only and does not extend to other rural or urban
healthcare facilities, limiting the generalizability of findings. The system excludes advanced
features such as laboratory integration, payment and billing, pharmacy management, patient
portals, and iOS compatibility. The evaluation covers only the initial prototype development and
early user feedback, and does not assess long-term system performance, cost-effectiveness, or
sustained implementation. External factors including internet instability, power outages, varying
digital literacy levels, and changes in healthcare policy are also outside the scope of this
study.

\subsection*{Significance of the Study}

This study benefits the following stakeholders in advancing rural healthcare delivery and digital
health innovation.

\textbf{Patients in Rural Communities.} Rural residents gain access to healthcare services
without costly travel, reducing time and financial burdens while ensuring continuity of care for
chronic conditions.

\textbf{Rural Healthcare Providers.} Providers benefit from automated records management,
streamlined appointment scheduling, and simplified specialist referrals, enabling more efficient
and accurate patient care.

\textbf{Community Health Workers.} Frontline workers gain faster access to patient information,
reducing paperwork and allowing greater focus on direct care, remote consultations, and early
detection of emerging health trends.

\textbf{Health Administrators and Local Government Officials.} Administrators gain data-driven
insights on service use, disease prevalence, and resource allocation, along with a replicable
model for digital health initiatives in other underserved areas.

\textbf{Technology Developers.} Developers gain practical, evidence-based insights into building
offline-capable, low-bandwidth systems designed for users with limited digital literacy in rural
settings.

\textbf{Academic Researchers.} The study provides validated frameworks and tools for designing
context-sensitive digital health solutions, filling existing literature gaps on integrated
telemedicine platforms in resource-limited environments.

\subsection*{Definition of Terms}

This section provides clear and precise definitions of key terms used throughout the study, both
operationally and conceptually.

\textbf{E-health.} E-health refers to the use of information and communication technologies in
healthcare delivery, including telemedicine, electronic medical records, and digital health
platforms \cite{WHO2023}. In this study, it encompasses all electronically delivered health
services through the proposed telemedicine platform for rural communities.

\textbf{E-health for All.} E-health for All is a principle that advocates for equal access to
digital health technologies regardless of location, income, or technical ability \cite{Kepper2024}.
In this study, it refers to the commitment to developing inclusive, user-friendly telemedicine
solutions specifically for rural communities in Naga City.

\textbf{Telemedicine.} Telemedicine is the delivery of healthcare services and clinical
information through communication technologies to overcome geographical barriers \cite{Li2025}.
In this study, it refers to the proposed remote healthcare delivery system developed for Naga
City.

\textbf{User-Centered Design.} User-centered design is a development methodology that focuses on
the needs, preferences, and limitations of end-users throughout the design process \cite{ISO9241}.
In this study, it refers to the participatory design approach used to create the telemedicine
platform for rural patients and healthcare providers in Naga City.

\textbf{Rural Healthcare.} Rural healthcare refers to the provision of medical services in
isolated or underserved areas with limited infrastructure and specialist access \cite{CDC2025}. In
this study, it refers to healthcare delivery in the barangays of Naga City outside the urban
center facing barriers such as long travel distances and limited specialist consultations.

\textbf{Telehealth Platform.} A telehealth platform is a digital health ecosystem integrating
provider directories, patient information systems, and scheduling tools to facilitate remote
healthcare delivery \cite{SPro2024}. In this study, it refers to the web-based system developed
for Naga City that integrates a doctors' database, electronic medical records, and appointment
scheduling into a single interface.

\textbf{Healthcare Accessibility.} Healthcare accessibility refers to the ease with which
individuals can obtain necessary health services considering physical, geographical, economic, and
organizational factors. In this study, it refers to how rural residents of Naga City can access
medical consultations and health information through the telemedicine platform, removing
geographical and time barriers to care.

%% ═══════════════════════════════════════════════════════════════════════════
%%  SECTION 2 — METHODOLOGIES
%% ═══════════════════════════════════════════════════════════════════════════
\section{Methodologies}

This section outlines the research methodology. It covers the research design, participants,
location, data collection methods and tools, ethical considerations, statistical analysis, project
timeline, and validation methods to ensure the reliability of the findings.

\subsection*{Research Design}

This study adopted a developmental-descriptive research design to document current healthcare
practices, identify operational challenges, and develop a telemedicine solution suited to the
needs of Barangay Panicuason RHU. Data were collected through surveys, interviews, observations,
and document analysis, following four phases: preparation, assessment, data collection, and
analysis.

\subsection*{Population and Sample}

The respondents were the 13 health facility personnel of Barangay Panicuason Rural Health Unit,
selected through total enumeration as primary end-users of the proposed system. The respondents
comprised 7 Dental Health Workers, 1 Midwife, 1 Encoder, 1 Barangay Nutrition Scholar (BNS), 1
Supplemental Feeding Volunteer (SFV), 1 Community-Based Rehabilitation Specialist (CBRS), and 1
Barangay Social Welfare Aide (BSWA).

For the succeeding phase involving rural community residents as secondary respondents, purposive
sampling will be applied, limited to residents with direct experience accessing the barangay
health facility. The required sample size will be determined using Slovin's formula:

\begin{equation}
  n = \frac{N}{1 + Ne^{2}}
  \label{eq:slovin}
\end{equation}

\noindent where:
\begin{itemize}[leftmargin=1in, label={}]
  \item $n$ = required sample size
  \item $N$ = total population of Barangay Panicuason
  \item $e$ = margin of error (set at 0.05 or 5\%)
\end{itemize}

This formula helps ensure our sample size is adequate to represent the barangay population while
maintaining statistical reliability.

\subsection*{Research Locale}

The study was conducted at Barangay Panicuason Rural Health Unit in Naga City, Camarines Sur, a
barangay situated outside the urban center and characterized by limited specialist availability,
long travel distances to medical facilities, and insufficient healthcare infrastructure. The
facility served as the primary healthcare provider for residents of Barangay Panicuason, offering
basic services including primary care, maternal and child health, immunization, and disease
prevention, and referring cases requiring specialized care to hospitals in Naga City proper.

\subsection*{Data Gathering Procedure}

This section describes the methods and tools the researchers used to capture data from
stakeholders essential to this study.

\textbf{Preparation Phase.} The researchers secured formal permissions from barangay officials
and local health officers, prepared consent forms, and coordinated with community leaders to
schedule observation and interview dates within the first two weeks.

\textbf{Document Analysis.} The researchers examined health records, referral systems, and
community health profiles from barangay health stations to document current healthcare access
patterns and identify existing challenges.

\textbf{Direct Observation.} The researchers observed and recorded how the health facility
personnel managed patient consultations, handled records, and coordinated scheduling and
appointments in their actual work setting.

\textbf{Surveys.} A structured questionnaire was administered to all 13 health facility personnel
to gather data on current practices, encountered challenges, importance of proposed system
features, and level of system acceptability.

\textbf{Interviews.} One-on-one interviews lasting 30 to 45 minutes were conducted with the 13
health facility personnel of Barangay Panicuason Rural Health Unit to explore current healthcare
practices, operational challenges, and readiness for a telemedicine solution.

\textbf{Data Analysis and Documentation.} Survey responses were tallied and computed using
weighted mean and frequency distribution, while interview data were organized by themes and used
as the basis for the final findings and recommendations.

\subsection*{Research Instrument}

This study used various research instruments to collect accurate and relevant data from
participants. These instruments are designed to gather both qualitative and quantitative
information to understand the current healthcare access situation and identify what is needed for
an effective telemedicine platform.

\textbf{Interview Guides.} Researchers prepared interview guides for rural residents of Barangay
Panicuason covering healthcare experiences, travel distance to facilities, medical costs including
transportation, common health problems, mobile and internet access, technology comfort level, and
desired telemedicine features. For health center staff, guides covered record-keeping systems,
data management challenges, internet connectivity, available devices, and infrastructure needs.
All guides included follow-up questions to explore key points further.

\textbf{Survey Questionnaires.} For healthcare providers and health workers, the questionnaire
covers demographic profile, years of experience, and current practices in patient records
management, doctor scheduling, and appointment systems using a 5-point frequency scale. It also
addresses challenges encountered in each of these areas, the importance of proposed system
features such as digital record storage, automated scheduling, online appointment booking, and
role-based access control, and the perceived level of system acceptability in terms of usability,
reliability, and data security.

\subsection*{Ethical Considerations}

This research followed strict ethical guidelines to protect all participants and ensure
responsible handling of operational data. The following measures were implemented:

\textbf{Voluntary Participation.} Participation was completely voluntary, with no pressure or
undue incentives influencing the decision to join. Participants may skip any question or withdraw
at any time without consequences or effect on their access to healthcare services.

\textbf{Informed Consent.} All participants received detailed information covering the study's
purpose, data collection and use, potential risks and benefits, confidentiality measures, and the
right to withdraw before agreeing to participate. Those who agreed signed an informed consent
form.

\textbf{Parental/Guardian Consent and Minor Assent.} For participants under 18, written parental
or guardian consent will be obtained prior to any research involvement, and minors aged 7--17
will provide written assent after receiving age-appropriate study information. Minors may decline
participation even if parents have given consent.

\textbf{Consent from Barangay Health Workers (BHWs).} BHWs signed informed consent forms
acknowledging permission to collect and use their personal information, professional insights, and
community health observations. The form clarified how their contributions will be used and ensured
their information is protected under data privacy regulations.

\textbf{Privacy and Confidentiality.} Individual identities were not revealed without explicit
written permission, and all interview recordings and transcripts were stored in password-protected
files accessible only to the research team. Researchers will not photograph or record any actual
patient medical records containing private health information.

\textbf{Data Protection.} All research data will be handled in accordance with the Data Privacy
Act of 2012 (Republic Act No.\ 10173), with personal identifiers removed and survey responses
kept anonymous. All data will be securely destroyed after thesis completion and any required
retention period.

\textbf{Minimizing Risk.} Research procedures present minimal risk, and researchers traveled to
participants' communities to reduce inconvenience. If participants felt uncomfortable during
interviews, sessions were paused or ended immediately with appropriate referrals provided if
needed.

\textbf{Cultural Sensitivity.} Researchers approached all interactions with cultural humility and
respect for local practices, collaborating with community leaders to ensure research activities
were culturally appropriate and scheduled at convenient times.

\textbf{Researcher Responsibilities.} The research team is committed to accurately recording and
reporting all data without fabrication or selective omission, properly citing all sources, and
honestly reporting the research process including challenges and limitations. Findings will be
shared with participating communities upon request.

\subsection*{Statistical Tools}

This section explains how we analyzed and interpreted the collected data to answer the research
questions.

\textbf{Descriptive Statistics.} The study employed descriptive statistical methods to summarize
and describe the characteristics of the sample and their responses. The following descriptive
statistics were used: percentage, frequency distribution, and weighted mean.

\textbf{Percentage.} To show the proportion of respondents who chose each answer option, we
calculated percentages using the formula:

\begin{equation}
  \% = \frac{f}{N} \times 100
  \label{eq:percentage}
\end{equation}

\noindent where $f$ = frequency or number of responses for a specific option, and $N$ = total
number of respondents. For example, if 45 out of 60 rural residents say they want mobile access
to health records, the percentage would be calculated as $(45/60) \times 100 = 75\%$.

\textbf{Frequency Distribution.} Frequency tables were created displaying how many respondents
selected each response option. For instance, for a satisfaction question with a 5-point scale,
the frequency distribution may be presented as shown in Table~\ref{tab:freq}.

\begin{table}[H]
  \centering
  \caption{Frequency Distribution}
  \label{tab:freq}
  \begin{tabular}{lcc}
    \toprule
    \textbf{Response Category} & \textbf{Number of Respondents} & \textbf{Percentage} \\
    \midrule
    Very Satisfied    & 5  & 8.3\%  \\
    Satisfied         & 25 & 41.7\% \\
    Neutral           & 15 & 25.0\% \\
    Dissatisfied      & 10 & 16.7\% \\
    Very Dissatisfied & 5  & 8.3\%  \\
    \midrule
    \textbf{Total}    & \textbf{60} & \textbf{100\%} \\
    \bottomrule
  \end{tabular}
\end{table}

\textbf{Weighted Mean.} For Likert-scale items, the weighted mean is calculated using the formula:

\begin{equation}
  \overline{x} = \frac{\sum (W \cdot f)}{N}
  \label{eq:wm}
\end{equation}

\noindent where $W$ = weight or scale value (1--5 for a 5-point Likert scale), $f$ = frequency or
number of respondents who chose that value, $N$ = total number of respondents, and $\Sigma$ =
summation symbol.

\subsection*{Statistical Treatment of Data}

The data gathered were interpreted using a five-point Likert scale to measure the respondents'
attitudes, opinions, perceptions, and preferences regarding the user-centered telemedicine system.

\textbf{Interpretation Scale.} To interpret the weighted mean values obtained from Likert-scale
questions, the following scale was used (Table~\ref{tab:likert}).

\begin{table}[H]
  \centering
  \caption{Likert Scale for Agreement on Telemedicine System}
  \label{tab:likert}
  \begin{tabular}{ccl}
    \toprule
    \textbf{Scale} & \textbf{Range} & \textbf{Interpretation} \\
    \midrule
    5 & 4.21--5.00 & Strongly Agree   \\
    4 & 3.41--4.20 & Agree            \\
    3 & 2.61--3.40 & Neutral          \\
    2 & 1.81--2.60 & Disagree         \\
    1 & 1.00--1.80 & Strongly Disagree \\
    \bottomrule
  \end{tabular}
\end{table}

\subsection*{System Design}

The system design process employed several modeling techniques such as use case and system
workflows to map user interactions and the system architecture. These techniques directly captured
the functional requirements, user actions, and step-by-step procedures within the system. This
ensured a thorough understanding for designing a robust and efficient telemedicine platform.

\textbf{Use Case.} The use case diagram in the study showed the specific interactions between
patients, healthcare providers, and system administrators with the system. It defined the actions
between these actors and the use cases, along with the functional requirements.

In the use case diagram, there were three actors: patients, doctors, and system administrators.
All actors can register and log in to the system. Patients and doctors utilize a mobile
application, while system administrators access a dedicated website application.

\begin{figure}[H]
  \centering
  % \includegraphics[width=0.85\textwidth]{figures/use_case.png}
  \fbox{\parbox{0.85\textwidth}{\centering\vspace{2cm}
    \textit{[Figure 3: Use Case Diagram --- insert image here]}
    \vspace{2cm}}}
  \caption{Use Case Diagram}
  \label{fig:usecase}
\end{figure}

\textbf{System Workflows.} The system workflows illustrate the sequential processes and decision
points for key operations within the E-Health for All Telemedicine System. These workflows ensure
standardized procedures and guide users through various system functionalities.

\begin{figure}[H]
  \centering
  % \includegraphics[width=0.85\textwidth]{figures/workflow.png}
  \fbox{\parbox{0.85\textwidth}{\centering\vspace{2cm}
    \textit{[Figure 4: System Workflow Diagram --- insert image here]}
    \vspace{2cm}}}
  \caption{System Workflow Diagram}
  \label{fig:workflow}
\end{figure}

\textbf{System Architecture.} The E-Health for All Telemedicine System architecture, shown in
Figure~\ref{fig:architecture}, connects three stakeholder groups---healthcare providers, rural
residents, and system administrators---through a layered structure spanning user devices,
connectivity, application services, integration, and data storage. The system is built around
three core components in the Application Layer: (1) a web-based platform for healthcare providers,
offering secure login, patient records management, clinical documentation, and digital prescription
and certificate generation; (2) a mobile application for rural residents, enabling appointment
scheduling, medical record access, video consultations, SMS reminders, health education content,
and emergency contacts; and (3) a doctor and health worker scheduling system, maintaining provider
profiles including specializations, availability, and consultation details.

\begin{figure}[H]
  \centering
  % \includegraphics[width=0.85\textwidth]{figures/architecture.png}
  \fbox{\parbox{0.85\textwidth}{\centering\vspace{2cm}
    \textit{[Figure 5: E-Health for All System Design Architecture --- insert image here]}
    \vspace{2cm}}}
  \caption{E-Health for All System Design Architecture}
  \label{fig:architecture}
\end{figure}

\subsection*{Validation and Testing Methods}

To ensure comprehensive validation and testing of system quality and reliability:

\textbf{Functionality Testing.} Unit testing verified individual components using test cases
covering normal inputs, edge cases, and invalid inputs. Integration testing confirmed that data
flowed correctly between modules and that integrated components worked together as expected.

\textbf{User Acceptance Testing.} It involved healthcare providers evaluating whether the system
met their needs through a checklist validating that all requirements were implemented.

\textbf{Performance Testing.} Load testing simulated 10, 20, and 30 concurrent users to ensure
acceptable performance under typical and peak loads. Response time was tracked with a target of
under 3 seconds for most operations, while offline functionality testing confirmed data entry and
retrieval without internet connection.

\textbf{Security Testing.} Access control verification confirmed users could only access
role-appropriate features and that unauthorized access attempts were blocked. Data encryption
validation confirmed sensitive information was encrypted in transit via HTTPS and at rest in the
database, with passwords hashed using strong algorithms and sessions tested for appropriate
expiration and hijacking prevention.

\textbf{Usability Testing.} The think-aloud protocol was conducted with 5 to 8 representative
users who verbalized thoughts while completing tasks as researchers observed and recorded
difficulties and errors. Metrics included task completion rates, time on-task, error rates, and
post-task satisfaction rating.

%% ═══════════════════════════════════════════════════════════════════════════
%%  SECTION 3 — RESULTS AND DISCUSSION (Pre-Oral)
%% ═══════════════════════════════════════════════════════════════════════════
\section{Results and Discussion}

This chapter presents the findings gathered from the survey administered to the health facility
personnel of Barangay Panicuason Rural Health Unit in Naga City. Results are organized according
to the specific research objectives: (1) current healthcare practices, (2) challenges encountered
in the current system, (3) system design plans as the prototype response to identified challenges,
and (4) the level of acceptability of the initial prototype in terms of usability, reliability,
and data security. Data are presented using weighted mean scores and frequency distributions,
discussed in relation to the reviewed literature, and linked to the theoretical frameworks guiding
the study.

\subsection*{Profile of Respondents}

A total of 13 health facility personnel completed the survey questionnaire during the data-gathering
phase. Table~\ref{tab:demographics} presents the demographic profile of the respondents.

\begin{table}[H]
  \centering
  \caption{Demographic Profile of Survey Respondents ($n = 13$)}
  \label{tab:demographics}
  \begin{tabular}{lcc}
    \toprule
    \textbf{Category} & \textbf{Frequency ($f$)} & \textbf{Percentage (\%)} \\
    \midrule
    \multicolumn{3}{l}{\textbf{A. Sex}} \\
    Female & 12 & 92.3\% \\
    Male   & 1  & 7.7\%  \\
    Total  & 13 & 100\%  \\
    \midrule
    \multicolumn{3}{l}{\textbf{B. Position / Designation}} \\
    Dental Health Worker (DHW)                      & 7 & 53.8\% \\
    Midwife                                          & 1 & 7.7\%  \\
    Encoder                                          & 1 & 7.7\%  \\
    Barangay Nutrition Scholar (BNS)                 & 1 & 7.7\%  \\
    Community Based Rehabilitation Specialist (CBRS) & 1 & 7.7\%  \\
    Barangay Social Welfare Aide (BSWA)              & 1 & 7.7\%  \\
    Total                                            & 13 & 100\% \\
    \midrule
    \multicolumn{3}{l}{\textbf{C. Age Bracket}} \\
    Below 25 years old    & 1  & 7.7\%  \\
    25--34 years old      & 6  & 46.2\% \\
    35--44 years old      & 4  & 30.8\% \\
    45 years old and above & 2 & 15.4\% \\
    Total                 & 13 & 100\%  \\
    \midrule
    \multicolumn{3}{l}{\textbf{D. Years of Experience}} \\
    Less than 1 year  & 0  & 0\%     \\
    1--3 years        & 1  & 7.7\%   \\
    4--6 years        & 3  & 30.77\% \\
    7--9 years        & 0  & 0\%     \\
    10 years and above & 8 & 61.54\% \\
    Total             & 13 & 100\%   \\
    \bottomrule
  \end{tabular}
\end{table}

The respondents were predominantly female (92.3\%), consistent with the largely female composition
of community health workforces in Panicuason. The majority were aged 25--34 (46.2\%), indicating
a relatively young and potentially technology-receptive group. In terms of work experience, most
respondents had 10 years or more of service (61.54\%), followed by those with 4--6 years
(30.77\%), while only a small proportion had 1--3 years (7.69\%) and none fell within the 7--9
year range. This suggests that the respondents generally possess substantial professional
experience, enabling them to provide informed and reliable assessments of current operational
practices. These demographics align with the Task--Technology Fit (TTF) theory's premise that
technology adoption is influenced by users' experiential familiarity with their work tasks
\cite{TTF_Newcastle}.

\subsection*{Current Practices of Rural Healthcare Providers}

Research Question 1 asked: ``What are the current practices of rural healthcare providers with
respect to patient records management, doctor and health worker scheduling, and appointment
systems?'' Table~\ref{tab:practices} presents the weighted mean scores for each indicator.

\textit{Scale: 5 = Always (4.21--5.00), 4 = Often (3.41--4.20), 3 = Sometimes (2.61--3.40),
2 = Rarely (1.81--2.60), 1 = Never (1.00--1.80)}

\begin{table}[H]
  \centering
  \caption{Current Healthcare Practices --- Weighted Mean Scores ($n = 13$)}
  \label{tab:practices}
  \begin{tabular}{p{9cm}cc}
    \toprule
    \textbf{Indicator} & \textbf{WM} & \textbf{Interpretation} \\
    \midrule
    \multicolumn{3}{l}{\textbf{A. Patient Records Management}} \\
    2.1 Records documented immediately after each consultation & 4.70 & Always \\
    2.2 Records stored in an organized and accessible manner   & 4.60 & Always \\
    2.3 Records updated regularly to reflect current health status & 4.50 & Always \\
    2.4 Records retrieved quickly and accurately when needed   & 4.60 & Always \\
    \textbf{Overall WM --- Patient Records Management}        & \textbf{4.60} & \textbf{Always} \\
    \midrule
    \multicolumn{3}{l}{\textbf{B. Doctor and Health Worker Scheduling}} \\
    2.5 Schedules planned and communicated in advance          & 4.20 & Often \\
    2.6 Availability clearly recorded and accessible to staff  & 4.20 & Often \\
    2.7 Schedule changes updated and communicated promptly     & 4.20 & Often \\
    2.8 Backup/replacement arranged when scheduled worker is unavailable & 4.50 & Always \\
    \textbf{Overall WM --- Doctor Scheduling}                  & \textbf{4.28} & \textbf{Often} \\
    \midrule
    \multicolumn{3}{l}{\textbf{C. Appointment Systems}} \\
    2.9 Appointments arranged in an orderly and systematic manner & 4.50 & Always \\
    2.10 Patients informed of appointment schedules in advance  & 4.50 & Always \\
    2.11 Appointment records maintained accurately to avoid double-booking & 4.30 & Often \\
    2.12 Patients attended within a reasonable waiting time     & 4.20 & Often \\
    \textbf{Overall WM --- Appointment Systems}                & \textbf{4.38} & \textbf{Often} \\
    \bottomrule
  \end{tabular}
\end{table}

The results revealed consistently high practice scores across all three domains. Patient records
management registered the highest overall weighted mean of 4.60 (Always), indicating that health
facility personnel demonstrate strong commitment to documenting, organizing, and retrieving patient
records. Notably, records are documented immediately after consultations (WM = 4.70), reflecting
adherence to standard clinical documentation protocols.

Doctor and health worker scheduling obtained an overall WM of 4.28 (Often), with backup
arrangement scoring highest at 4.50 (Always), suggesting that contingency planning is
well-established at the facility. Appointment systems averaged 4.38 (Often), with systematic
arrangement and advance patient notification both scoring 4.50 (Always). These high practice
scores align with findings by Santos, Mendoza, and Cruz \cite{Santos2021}, who noted that
Philippine rural health units have established basic operational protocols, though implementation
quality varies significantly between facilities.

These findings support the Task--Technology Fit theory \cite{TTF_Newcastle} in that the existing
manual practices are well-established and purposive. A digital system that replicates and enhances
these existing workflows rather than replacing them is more likely to achieve strong adoption
among health workers who are already familiar with structured documentation and scheduling
routines.

\subsection*{Challenges Encountered in the Current System}

Research Question 2 asked: ``What are the challenges encountered in the current system?''
Table~\ref{tab:challenges} presents the weighted mean scores for each challenge indicator.

\textit{Scale: 5 = Always (4.21--5.00), 4 = Often (3.41--4.20), 3 = Sometimes (2.61--3.40),
2 = Rarely (1.81--2.60), 1 = Never (1.00--1.80)}

\begin{table}[H]
  \centering
  \caption{Challenges in Current Healthcare Practices --- Weighted Mean Scores ($n = 13$)}
  \label{tab:challenges}
  \begin{tabular}{p{9cm}cc}
    \toprule
    \textbf{Challenge Indicator} & \textbf{WM} & \textbf{Interpretation} \\
    \midrule
    \multicolumn{3}{l}{\textbf{A. Patient Records Management}} \\
    3.1 Patient records are difficult to organize and manage        & 3.20 & Sometimes \\
    3.2 Patient records are missing or incomplete                   & 2.80 & Sometimes \\
    3.3 Retrieving a patient's record takes too much time           & 3.00 & Sometimes \\
    3.4 Records are lost or damaged due to current storage method   & 2.80 & Sometimes \\
    \textbf{Overall WM --- Patient Records Challenges}              & \textbf{2.95} & \textbf{Sometimes} \\
    \midrule
    \multicolumn{3}{l}{\textbf{B. Doctor and Health Worker Scheduling}} \\
    3.5 Scheduling conflicts among doctors/health workers occur     & 3.00 & Sometimes \\
    3.6 No clear or updated record of availability                  & 2.90 & Sometimes \\
    3.7 Health workers assigned to multiple duties simultaneously   & 2.70 & Sometimes \\
    3.8 Schedule changes not communicated to concerned staff        & 2.70 & Sometimes \\
    \textbf{Overall WM --- Scheduling Challenges}                   & \textbf{2.83} & \textbf{Sometimes} \\
    \midrule
    \multicolumn{3}{l}{\textbf{C. Appointment Systems}} \\
    3.9 Patients experience long waiting times before being attended & 3.00 & Sometimes \\
    3.10 Delays in processing patient appointments                  & 2.90 & Sometimes \\
    3.11 Double-bookings or missed appointments occur               & 2.70 & Sometimes \\
    3.12 No organized system for tracking or confirming appointments & 3.00 & Sometimes \\
    \textbf{Overall WM --- Appointment Challenges}                  & \textbf{2.90} & \textbf{Sometimes} \\
    \bottomrule
  \end{tabular}
\end{table}

All challenge indicators were rated in the ``Sometimes'' range (WM = 2.61--3.40), confirming
that while problems were not constant, they occurred with enough regularity to warrant systematic
intervention. The most frequently cited challenge across domains was the difficulty in organizing
and managing patient records (WM = 3.20), consistent with findings by D'Amore et al.\
\cite{DAmore2021} that paper-based record systems in rural health facilities are particularly
susceptible to organizational lapses due to high patient volumes and limited administrative
staffing.

Scheduling challenges averaged 2.83, with scheduling conflicts (WM = 3.00) and unclear
availability records (WM = 2.90) as primary concerns. This aligns with Mendoza, Cruz, and Santos
\cite{Mendoza2023}, who found that the absence of a centralized scheduling system in Philippine
rural health units contributes to coordination failures, particularly during peak consultation
periods. Appointment system challenges averaged 2.90, with the absence of a tracking system
(WM = 3.00) and long patient waiting times (WM = 3.00) as the most persistent issues. Del
Rosario \cite{DelRosario2024} similarly found that the lack of digital appointment management in
Philippine rural health centers contributes significantly to extended waiting times and patient
dissatisfaction.

These ``Sometimes'' ratings, while not alarming individually, collectively defined a pattern of
recurring inefficiency that compounds over time. They validated the need for a digital solution
that addresses record organization, scheduling conflict detection, appointment tracking, and
automated reminders---all of which are core features of the proposed E-Health for All system.

\subsection*{System Design and Plans for the Prototype}

Research Question 3 asked: ``What system can be designed to address the identified challenges?''
This section presents the importance ratings for proposed system features.

\textit{Scale: 5 = Very Important (4.21--5.00), 4 = Important (3.41--4.20), 3 = Moderately
Important (2.61--3.40), 2 = Slightly Important (1.81--2.60), 1 = Not Important (1.00--1.80)}

\begin{table}[H]
  \centering
  \caption{Importance of Proposed Features --- Weighted Mean Scores ($n = 13$)}
  \label{tab:features}
  \begin{tabular}{p{9cm}cc}
    \toprule
    \textbf{Proposed System Feature} & \textbf{WM} & \textbf{Interpretation} \\
    \midrule
    4.1 Digital storage and easy retrieval of patient records     & 4.50 & Very Important \\
    4.2 Automated doctor/worker scheduling with conflict detection & 4.50 & Very Important \\
    4.3 Online patient appointment booking                         & 4.60 & Very Important \\
    4.4 Automated appointment reminders and notifications          & 4.80 & Very Important \\
    4.5 Secure login and role-based access control                 & 4.90 & Very Important \\
    \midrule
    \textbf{Overall Weighted Mean}                                 & \textbf{4.66} & \textbf{Very Important} \\
    \bottomrule
  \end{tabular}
\end{table}

All five proposed features received ``Very Important'' ratings, with an overall weighted mean of
4.66. Secure login and role-based access control (WM = 4.90) ranked highest, reflecting a strong
awareness among health workers of data privacy obligations under the Data Privacy Act of 2012.
Automated reminders (WM = 4.80) ranked second, directly addressing the challenge of missed
appointments (Item 3.11, WM = 2.70) and long waiting times (Item 3.9, WM = 3.00). Online
appointment booking (WM = 4.60) and digital storage (WM = 4.50) also scored highly, validating
the core functional design of the proposed system. These ratings are consistent with the TAM
framework \cite{TAM_Newcastle}, where perceived usefulness---as evidenced by the high importance
ratings---is a primary determinant of technology adoption intention.

%% ═══════════════════════════════════════════════════════════════════════════
%%  SECTION 4 — CONCLUSION AND RECOMMENDATIONS
%% ═══════════════════════════════════════════════════════════════════════════
\section{Conclusion and Recommendations}

This section provides a concise synthesis of the study's key findings and reflects on their
broader significance in advancing rural healthcare delivery through digital innovation.

This study aimed to design and develop E-Health for All, a user-centered telemedicine platform for
Barangay Panicuason Rural Health Unit in Naga City, addressing inefficiencies in patient records
management, health worker scheduling, and appointment systems. Using a developmental-descriptive
research design, data were gathered from 13 health facility personnel to document current
practices, identify challenges, evaluate proposed system features, and assess prototype
acceptability.

The findings reveal that health facility personnel had already established strong baseline
practices across all three operational domains, indicating a structured and disciplined workforce.
This is significant because it confirms that the E-Health for All system enters a context where
existing workflows are well-established, making adoption more likely when the platform is designed
to enhance rather than replace familiar routines---consistent with the Task--Technology Fit (TTF)
theory.

Despite these strong practices, recurring challenges were confirmed across all domains.
Difficulties in organizing physical records, the absence of a systematic appointment tracking
mechanism, and scheduling conflicts among health workers were the most consistently reported
concerns. These findings collectively validate the study's core rationale that an integrated
telemedicine platform is both necessary and timely for a facility of this scale and context.

All proposed system features were rated very important, with data security, automated reminders,
and online booking receiving the highest endorsements---directly corresponding to the identified
challenges. This alignment confirms that the system is grounded in users' actual needs, consistent
with the Technology Acceptance Model (TAM). The respondent profile---predominantly female,
relatively young, and experienced---further supports readiness for a managed digital transition
as affirmed by the Diffusion of Innovations (DOI) theory.

Overall, this study demonstrates that E-Health for All is both technically feasible and
operationally justified for Barangay Panicuason RHU. Beyond its immediate application, it offers
a replicable framework for designing context-sensitive telemedicine solutions in similarly
resource-limited settings.

\subsection*{Recommendations}

Based on the conclusions of this study, the following recommendations are directed to the key
stakeholders involved in the deployment, improvement, and continued research of the E-Health for
All telemedicine system.

\textbf{For the Barangay Panicuason Health Unit and the Local Government Unit.} The local
government unit should pursue the formal adoption and full deployment of the E-Health for All
system by allocating a dedicated budget for device procurement, internet connectivity, and system
maintenance. A structured onboarding program should be conducted prior to deployment to ensure a
smooth transition from paper-based to digital operations.

\textbf{For System Developers and Future Development Teams.} Developers should expand the system
in future iterations to include laboratory result integration, pharmacy management, a
patient-facing portal, and iOS compatibility. Offline functionality should be strengthened to
maintain data access during internet outages, and all development must strictly uphold the Data
Privacy Act of 2012 through encryption, role-based access controls, and regular security auditing.

\textbf{For Health Facility Personnel.} Health workers should actively engage during the pilot
phase by documenting usability issues and communicating improvement suggestions through formal
feedback channels. Their sustained adoption and cooperation will be essential to the system's
continuous refinement and long-term success.

\textbf{For Academic Researchers and Future Studies.} Future researchers should extend this study
by including community residents as secondary respondents and conducting longitudinal evaluations
to assess long-term system impact. Comparative studies across multiple rural health units are
recommended, and future instruments should incorporate the System Usability Scale (SUS) and
Cronbach's Alpha to strengthen methodological rigor.

\textbf{For Policymakers and the Department of Health.} The Department of Health should consider
the E-Health for All model as a pilot reference for rural health unit digitalization nationwide.
Policymakers should provide technical and financial support to local government units pursuing
similar initiatives and establish interoperability standards to connect locally developed platforms
with national health systems such as iHOMIS and the National Health Data Repository.

%% ═══════════════════════════════════════════════════════════════════════════
%%  REFERENCES
%% ═══════════════════════════════════════════════════════════════════════════
\bibliographystyle{abbrvnat}
\bibliography{references}

\end{document}